\section{Purpose}

\textbf{[STIPULATION]} This document develops the complete field theory of FFGFT:
explicit mode functions on the 4-torus, the free Lagrangian, the bridge formula
between Lagrangian and mode operator, the Dirac and Schrödinger equations as
low-energy limits, the binding picture of masses, and the non-closure theorem.
It is the link between the geometric structure (A030/A040) and the sectors
(A100\,ff.) -- it shows that masses do not appear as parameters in the Lagrangian
but emerge as eigenvalues of the mode operator.

\section{The Basic Geometry of the Torus}

\textbf{[B]} The fundamental carrier is the compact 4-torus $T^4$ with the
characteristic length scale
\[
  L_T = \frac{\ell_P}{\xi}\approx7500\,\ell_P\approx6{.}14\times10^{-16}\,\text{GeV}^{-1}.
\]
(Not to be confused with the minimal length scale $L_0=\xi\cdot\ell_P\approx2{.}15\times10^{-39}\,$m
from A015.) This is the \emph{fundamental period} of the torus: all particles
are modes on $L_T$ -- like overtones of a vibrating string in three spatial
dimensions. The Compton wavelength $L_i=2\pi/m_i$ is the wavelength of the
respective mode, not the size of a separate torus.

\textbf{[B]} The FFGFT torus carries a fractal correction of the dimension:
\[
  D_f = 3-\xi\approx 2{.}9998667.
\]
This minimal deviation is the origin of all particle masses. It generates a
non-linear dispersion relation -- analogous to an anharmonic crystal lattice.
On an exactly three-dimensional torus, masses would be integer-proportional to
the wave number; the observed ratios $m_\mu/m_e\approx206{.}8$ and
$m_\tau/m_e\approx3477$ force the fractal dispersion.

\section{Quantum Numbers of the Modes}

\textbf{[B]} Each mode on the torus carries three quantum numbers: $n$ (principal
quantum number), $l$ (orbital angular momentum, transverse component), $j$ (spin
quantum number, $j=l\pm\tfrac12$ for fermions). The wave vector is quantised:
\[
  \mathbf k_{n,l} = \frac{2\pi}{L_0}(k_x,k_y,k_z),\qquad k_x,k_y,k_z\in\mathbb Z,
\]
with $n=k_x^2+k_y^2+k_z^2$ and $l=\sqrt{k_x^2+k_y^2}$. The assignment for
the leptons:

\begin{center}
\renewcommand{\arraystretch}{1.3}
{\small
\begin{tabular}{lccccc}
\toprule
Particle & $n$ & $l$ & $(k_x,k_y,k_z)$ & $r_i$ & $p_i$ \\
\midrule
Electron & 1 & 0 & $(1,0,0)$ & $4/3$ & $3/2$ \\
Muon     & 2 & 1 & $(1,1,0)$ & $16/5$ & $1$ \\
Tau      & 3 & 2 & $(1,1,1)$ & $25/9$ & $2/3$ \\
\bottomrule
\end{tabular}
}
\renewcommand{\arraystretch}{1.0}
\end{center}

\textbf{[K]} The lepton exponents $(p_e,p_\mu,p_\tau)=(3/2,1,2/3)$ form a
geometric sequence with ratio $2/3$ -- the inverse of the tetrahedral factor
$4/3$ from $\xi=4/3\times10^{-4}$:
\[
  p_{n+1} = \tfrac23\,p_n.
\]
Up-type quarks (up to Charm) follow the same scheme with a step skipped;
Down-type quarks an arithmetic sequence ($\Delta p=-1/2$); the top quark
($p_t=-1/3$, $r_t=1/28$) breaks the scheme -- no known geometric reason
(open, A150).

\section{Explicit Mode Functions}

\textbf{[B]} The eigenfunctions of the d'Alembertian $\Box=\partial_t^2-\nabla^2$
on $T^3\times\mathbb R$ are:
\[
  \boxed{\;
    \phi_{n,l,j}(t,\mathbf x)
    = \frac{1}{\sqrt V}\,
      e^{i\mathbf k_{n,l}\cdot\mathbf x}\cdot
      Y_{l,j}(\vartheta,\varphi)\cdot
      e^{-im_{n,l,j}\,t}
  \;}
\]
with torus volume $V=L_0^3$ and spherical harmonics $Y_{l,j}$. The modes are
orthonormal:
\[
  \int_{T^3}\phi_{n,l,j}^*(\mathbf x)\,\phi_{n',l',j'}(\mathbf x)\,d^3x
  = \delta_{nn'}\delta_{ll'}\delta_{jj'}
\]
and satisfy the Klein--Gordon equation:
\[
  \Box\,\phi_{n,l,j} = -m_{n,l,j}^2\,\phi_{n,l,j}.
\]
The spin structure follows from the winding numbers of the torus: spin-$\tfrac12$
corresponds to the winding $(n_\theta,n_\phi)=(1,1)$; the 720° symmetry follows
directly from the topology.

\section{The Lagrangian}

\textbf{[B]} The starting point is the free Klein--Gordon Lagrangian:
\[
  \mathcal L_0 = \tfrac12(\partial\,\delta m)^2,
\]
the scalar field $\delta m(x)$ describes mass fluctuations on the torus; the
equation of motion $\Box\,\delta m=0$ describes massless oscillations. The
fractal geometry adds a correction term generating the non-linear dispersion:
\[
  \Delta\mathcal L_{\text{Torus}}
  = \varepsilon\,\xi\!\left[
    \frac{f(\vartheta)}{2}(\partial\,\delta m)^2
    -k^{\mu\nu}(\vartheta)\,\partial_\mu\delta m\,\partial_\nu\delta m
  \right].
\]
The geometric functions $f(\vartheta)$ and $k^{\mu\nu}(\vartheta)$ encode the
curvature of the fractal torus. The complete Lagrangian:
\[
  \boxed{\;\mathcal L = \mathcal L_0 + \Delta\mathcal L_{\text{Torus}}\;}
\]

\textbf{[K]} \textbf{Connection to the Standard Model.} For $\xi\to0$,
$\Delta\mathcal L_{\text{Torus}}$ vanishes; FFGFT contains the SM Lagrangian as
a limiting case. The difference from Yukawa: in the SM the coupling $y_i$ is a
free parameter; in FFGFT the corresponding term $\xi\cdot m_\ell$ is
geometrically fixed.

\section{The Mode Operator}

\textbf{[B]} The Hilbert space of torus modes:
\[
  \mathcal H = \bigoplus_{n,l,j}\mathbb C\cdot\mathbf e_{n,l,j},
  \qquad
  \langle\mathbf e_{n,l,j},\mathbf e_{n',l',j'}\rangle = \delta_{nn'}\delta_{ll'}\delta_{jj'}.
\]
The field operator:
\[
  \Phi(x) = \sum_{n,l,j} m_{n,l,j}\cdot\phi_{n,l,j}(x)\cdot P_{n,l,j},
\]
with projectors $P_i=|\mathbf e_i\rangle\langle\mathbf e_i|$. Since
$P_iP_j=\delta_{ij}P_i$, we have $\Phi^2=\sum_i m_i^2 P_i$.

\section{The Bridge Formula: Lagrangian and Mode Operator}

\textbf{[B]} The core of the field theory: the Lagrangian and the mode operator
describe the same physics. The bridge formula in six steps:

\[
  \boxed{\;
    \int\tfrac12(\partial\,\delta m)^2\,d^4x
    = \tfrac12\langle\psi|\Phi^2|\psi\rangle
  \;}
\]

\textbf{Step 1 -- Mode expansion:}
$\delta m(x) = \sum_i a_i\phi_i(x)$.

\textbf{Step 2 -- Substitution:}
$\int\mathcal L_0\,d^4x
= \tfrac12\sum_{i,j}a_ia_j\int(\partial\phi_i)(\partial\phi_j)\,d^4x$.

\textbf{Step 3 -- Orthogonality:}
$\int(\partial\phi_i)(\partial\phi_j)\,d^4x=0$ for $i\neq j$.
Hence $= \tfrac12\sum_ia_i^2\int(\partial\phi_i)^2\,d^4x$.

\textbf{Step 4 -- Klein--Gordon (partial integration):}
$\int(\partial\phi_i)^2\,d^4x = m_i^2$ (normalisation).

\textbf{Step 5 -- Summary:}
$\int\mathcal L_0\,d^4x = \tfrac12\sum_ia_i^2 m_i^2$.

\textbf{Step 6 -- Mode operator:}
$\langle\psi|\Phi^2|\psi\rangle = \sum_ia_i^2 m_i^2$.
Comparison yields the bridge formula.

\textbf{[B]} Three description levels are thereby equivalent:

\begin{center}
\renewcommand{\arraystretch}{1.3}
{\small
\begin{tabular}{ll}
\toprule
Level & Formula \\
\midrule
Lagrangian     & $\mathcal L_0 = \tfrac12(\partial\,\delta m)^2$ \\
Mode operator  & $\Phi = \sum_i m_i P_i$ \\
Bridge         & $\int\mathcal L_0\,d^4x = \tfrac12\langle\psi|\Phi^2|\psi\rangle$ \\
\bottomrule
\end{tabular}
}
\renewcommand{\arraystretch}{1.0}
\end{center}

Mass arises not as a Lagrangian parameter but as an eigenvalue of the mode
operator -- this is the structural difference from the Standard Model.

\section{Mass Generation as a Binding Phenomenon}

\textbf{[B]} The eigenvalue equation on the torus:
\[
  \bigl[\Box + \hat V(\xi,D_f)\bigr]\phi_{n,l,j} = -m_{n,l,j}^2\,\phi_{n,l,j}.
\]
The mass results as the difference of kinetic energy and potential:
\[
  m_i^2
  = \underbrace{\Bigl(\frac{2\pi n}{L_0}\Bigr)^2}_{\text{kinetic}}
  + \underbrace{\langle\phi_i|\hat V|\phi_i\rangle}_{\approx\,-(2\pi n/L_0)^2}.
\]

\textbf{[S]} The kinetic energy is enormous ($\sim10^{32}\,\text{GeV}^2$),
the potential almost exactly opposite. The mass is the tiny remainder. This is
not fine tuning but the defining property of bound states. \emph{Crystal
analogy}: one does not ask why a crystal produces a particular phonon
dispersion -- the dispersion curve \emph{is} the crystal structure in momentum
space. \emph{Hydrogen analogy}: kinetic energy $\approx13{.}6\,\text{eV}$,
Coulomb potential $\approx-27{.}2\,\text{eV}$, binding energy $-13{.}6\,\text{eV}$
-- no one asks why the potential compensates the kinetic energy so precisely,
that follows from the Schrödinger equation. In FFGFT: the particles are bound
torus modes; the potential $\hat V$ is uniquely determined by $\xi$ and $D_f$
-- no free parameter.

\section{The Non-Closure Theorem}

\textbf{[K]} For all pairs of distinct torus modes $i\neq j$:
\[
  \forall i\neq j:\; (r_i\cdot r_j,\; p_i+p_j)\notin\mathcal P,
\]
where $\mathcal P$ is the set of all pairs $(r,p)$ admitted in FFGFT.

Example: $r_e\cdot r_\mu = \tfrac43\cdot\tfrac{16}{5} = \tfrac{64}{15}$,
$p_e+p_\mu = \tfrac32+1 = \tfrac52$. In the complete mode table there is no
particle with $r=64/15$ and $p=5/2$.

\textbf{[B]} Why does the theorem hold? The cause lies in the orthogonality of
the modes: in Hilbert space the product of two distinct basis vectors is
structurally meaningless -- a category violation. $r_e\cdot r_\mu$ is computable
as a number, but belongs to no mode of the torus.

\textbf{[K]} Practical rule: \emph{Permitted} is numerical evaluation
($m_e\cdot m_\mu\approx53{.}9\,\text{MeV}^2$) and the mass formula
$m_i=r_i\cdot\xi^{p_i}\cdot v$ for fixed $i$. \emph{Forbidden} is symbolic
cross-multiplication: $(r_i\cdot r_j)\cdot\xi^{p_i+p_j}\cdot v$ for $i\neq j$.

\section{Dirac Line and Schrödinger Limit}

\textbf{[B]} Within the Lagrangian framework the geometric equation replaces the
full Dirac equation by:
\[
  \partial^2\,\Delta m = 0.
\]
The fundamental structure is the Clifford algebra
$\mathbf e_\mu\mathbf e_\nu+\mathbf e_\nu\mathbf e_\mu=2g_{\mu\nu}$; the
$4\times4$ $\gamma$-matrices are retained as a representation of the Clifford
generators for concrete computations -- they are not eliminated but traced back
to their geometric meaning. Spin manifests on the fractal torus as a topological
property (vacuum-phase winding). A relativistic state:
\[
  \Phi(x,t) = \rho(x,t)\,e^{i\theta(x,t)}\cdot\chi(\sigma),
\]
with spin component $\chi(\sigma)$ from the geometric knot modes.

\textbf{[B]} In the non-relativistic limit, the phase equation of the vacuum
field reduces to the standard Schrödinger equation:
\[
  i\hbar\,\frac{\partial\psi}{\partial t}
  = -\frac{\hbar^2}{2m}\nabla^2\psi + V\psi,
\]
with $\psi\propto\sqrt\rho\,e^{i\theta}$, effective mass $m$ from the knot
inertia, and potential $V$ from external $\rho$-perturbations. The Schrödinger
equation is \emph{not fundamental} but an effective equation for slow,
non-relativistic knot excitations -- the low-energy limit of the FFGFT mode
structure. This limit follows from the Lagrangian framework via the DVFT polar
decomposition $\Phi\to\sqrt\rho\,e^{i\theta}$.

\textbf{[B]} The transition Lagrangian $\to$ Schrödinger equation in three
stages: (1) mode superposition in the low-energy limit:
$\Phi(x)\xrightarrow{\text{low energy}}\sqrt{\rho(x,t)}\,e^{i\theta(x,t)}$
with energy density $\rho$ and coherent phase $\theta$. (2) The phase dynamics
of $\theta$ satisfies the Schrödinger equation in the non-relativistic limit.
(3) Masses follow from $k_t$-windings via $m_ic^2=\hbar f(k_t^{(i)})$; the
lepton exponents $(3/2,1,2/3)$ are the geometric sequence of this tower.

\section{Algebras, Recursion Operator, and Scale Structure}

\textbf{[B]} The recursion operator $\mathcal R$ formalises the fractal scaling
structure, diagonal on the mode space:
\[
  \mathcal R(\mathbf e_{n,l,j}) = \xi^{p(n,l,j)-1}\cdot K_\text{frak}\cdot\mathbf e_{n,l,j}.
\]
The time-mass duality is encoded as a field symmetry:
$\mathcal R:\;t\to\xi t,\;m\to\xi^{-1}m,\;\Phi\to\Phi$. The mode algebra for
free waves: $\mathbf e_{\mathbf k_i}\star\mathbf e_{\mathbf k_j}
=\mathbf e_{\mathbf k_i+\mathbf k_j}$, with structure constants
$C_{ij}^k=\delta_{k,i+j\bmod\text{torus period}}$.

\textbf{[K]} The scale running of the coupling -- in FFGFT-native language: an
$\mathcal R$-induced scale correction, not a renormalisation procedure -- gives
in the QED analogy:
\[
  \alpha^{-1}(\mu) = \alpha^{-1}(m_e)
  - \frac{\beta_0}{2\pi}\ln\frac\mu{m_e}
  - \frac{\beta_0\xi}{4\pi}\Bigl(\ln\frac\mu{m_e}\Bigr)^2.
\]
The $\xi$-term is the FFGFT-specific addition; $\xi$ is not a fixed point of
the $\beta$-function but a geometrically fixed constant (A020).

\section{Methodological Position of the Field Theory}

\textbf{[K]} \textbf{Reduction is not increased precision.} Where the Standard
Model is experimentally confirmed (QED scattering amplitudes to $\sim10^{-12}$),
FFGFT can only reproduce these plus $\xi$-corrections of order $10^{-4}$, not
qualitatively surpass them. The FFGFT mass formulas give values at $\sim1\,\%$
agreement with PDG values -- that is a \emph{structural proof} (masses follow
from geometry, not free parameters), not a precision calculation. The reduction
from $\sim25$ free SM parameters to a single geometric parameter $\xi$ is the
quantifiable statement.

\textbf{[S]} \textbf{Circularity of the comparison basis.} PDG values are not
model-free raw data; they contain assumptions about units, gauge conventions, and
SI-2019-reform definitions. An FFGFT calculation in natural units ($\hbar=c=1$)
avoids the accumulation of model assumptions; as soon as it is compared against
PDG data, it inherits their prior corrections. The $\sim1\,\%$ comparison is
therefore to be read in two directions: confirmation of the geometric structure
\emph{and} indication of the latitude of the circular measurement basis.

\section{Open Mathematical Tasks}

\textbf{The Hopf algebra for interactions is not worked out.} The free theory
has structure constants $C_{ij}^k=\delta_{k,i+j}$ (wave vector addition).
Complete Feynman rules from the structure constants and the non-trivial fixed
point of the $\beta$-function (whether $\mu^*$ exists and what it means
physically) are open.

\textbf{Top quark, Down-type step size, quark $r_i$ deep structure.} For leptons
and Up-type quarks (up to Charm) the geometric $2/3$-sequence of exponents is
known; for the Top ($p_t=-1/3$) and Down-type (arithmetic sequence $\Delta p=-1/2$)
the geometric justification is missing. The $\alpha$-$\pi$-$\sqrt3$ deep
structure of the lepton $r_i$ is partially present (A130); the analogue for
quark $r_i$ is missing (A150).

\textbf{The test condition is structural, not quantitative.} The test
$\langle\phi_i|\hat V|\phi_i\rangle = r_i^2\xi^{2p_i}v^2-(2\pi n/L_0)^2$ is the
bottleneck: explicitly deriving $\Delta\mathcal L_\text{Torus}$ from the torus
geometry and showing that the eigenvalues agree with the mass table -- as phonon
dispersion from lattice forces of a crystal.

\section{Verification Scripts}

\begin{itemize}[leftmargin=2em,itemsep=2pt,topsep=2pt]
  \item Mode functions and orthonormality: verification
    $\int\phi_i^*\phi_j\,d^3x=\delta_{ij}$ for $i,j\in\{e,\mu,\tau\}$
    numerically:
    \texttt{python/A\_Serie\_Skripte/a075\_modefunktionen.py}
  \item Bridge formula: $\int\tfrac12(\partial\delta m)^2\,d^4x
    = \tfrac12\langle\psi|\Phi^2|\psi\rangle$ for random coefficients $(a_i)$:
    \texttt{python/A\_Serie\_Skripte/a075\_bruecke.py}
  \item Non-closure theorem: verification that $(r_e r_\mu, p_e+p_\mu)$
    does not lie in $\mathcal P$:
    \texttt{python/A\_Serie\_Skripte/a075\_nichtabschluss.py}
\end{itemize}

\section{Symbol Glossary}

Symbols used throughout the entire A-series are explained in A010.
Specific to this document:

\begin{center}
\renewcommand{\arraystretch}{1.3}
{\small\begin{tabular}{lp{9cm}}
\toprule
Symbol & Meaning \\
\midrule
$\Box = \partial_t^2 - \nabla^2$ & d'Alembertian (wave operator) \\
$\mathcal{L}$ & Lagrangian density \\
$\phi_{n,l,j}(x)$ & Mode function on $T^4$; quantum numbers as in A100 \\
$m_{n,l,j}$ & Mass of mode $(n,l,j)$ \\
$\Psi_\ell$ & Lepton spinor field \\
$g_T^\ell = m_\ell\xi$ & Mass-proportional coupling of the time field \\
\bottomrule
\end{tabular}}
\renewcommand{\arraystretch}{1.0}
\end{center}

\section{Open Questions}

\textbf{The complete form of $\Delta\mathcal L_{\text{Torus}}$} -- the geometric
functions $f(\vartheta)$ and $k^{\mu\nu}(\vartheta)$ -- is not derived in
closed form from the torus geometry. The correction factor $K_\text{frak}
=1-100\xi\approx0{.}9867$ is known (A040), the complete tensor function is not.

\textbf{The quark-sector extension} of the mode table and the non-closure theorem
to all twelve fermions is worked out for leptons and Up-type quarks (up to Charm),
but not for Down-type and Top (A150).

\textbf{The Hopf algebra structure for interactions} and complete Feynman rules
are outstanding.

\section{Supersedes}

This document subsumes: Dok.~202 (complete FFGFT field theory: fundamental scale
$L_0$, quantum numbers of modes, explicit mode functions $\phi_{n,l,j}$, free
Lagrangian $\mathcal L_0$, torus correction term, mode operator $\Phi=\sum m_i
P_i$, bridge formula in six steps, non-closure theorem, binding picture of
masses, Dirac and Schrödinger limits, recursion operator $\mathcal R$, scale
structure). \emph{Not} taken over: the large-scale sector from Dok.~202 (DVFT
line on vacuum field effects at galactic scales) -- that does not belong in the
A-series core (A010); the section on circularity of the maximal-formula notation
is left to the correction programme.

\section{Use of AI Tools}

This document was created with the assistance of AI tools.
Responsibility for content and errors rests entirely with the author.
Full wording of the rules in A010.
